\documentclass[a4paper,12pt]{article}
\usepackage[utf8]{inputenc}
\usepackage[ngerman]{babel}
\usepackage{amsmath,amssymb}
\usepackage{geometry}
\geometry{margin=2.5cm}
\usepackage{tcolorbox}
\tcbuselibrary{breakable}

\newtcolorbox{merkbox}[1][]{colback=blue!5, colframe=blue!40, fonttitle=\bfseries, title={Merke}, breakable, #1}
\newtcolorbox{analogie}[1][]{colback=green!5, colframe=green!40, fonttitle=\bfseries, title={Analogie}, breakable, #1}

\title{Erkl\"arung -- HW03: Untervektorr\"aume}
\author{}
\date{}

\begin{document}
\maketitle

%=============================================================================
%  TEIL 1: GRUNDLAGEN
%=============================================================================
\part*{Teil 1: Grundlagen -- Was du wissen musst, BEVOR du die Aufgaben machst}

%-----------------------------------------------------------------------------
\section*{1. Was ist ein Vektorraum?}

Ein \textbf{Vektorraum} ist einfach eine Menge von ``Objekten'', mit denen man zwei Dinge machen kann:
\begin{itemize}
    \item \textbf{Addieren:} Zwei Objekte zusammenz\"ahlen $\rightarrow$ Ergebnis ist wieder so ein Objekt.
    \item \textbf{Skalarmultiplikation:} Ein Objekt mit einer reellen Zahl multiplizieren $\rightarrow$ Ergebnis ist wieder so ein Objekt.
\end{itemize}

\begin{analogie}
Stell dir einen \textbf{Club} vor. Der Club hat Regeln (die Vektorraum-Axiome). Jedes Mitglied ist ein ``Vektor''. Die Regeln sagen zum Beispiel: Wenn du zwei Mitglieder ``addierst'', kommt wieder ein Mitglied raus. Wenn du ein Mitglied ``skalierst'', kommt wieder ein Mitglied raus.
\end{analogie}

\textbf{Beispiele f\"ur Vektorr\"aume:}
\begin{itemize}
    \item $\mathbb{R}^n$ -- die ``normalen'' Vektoren, z.B.\ $(1, 2, 3) \in \mathbb{R}^3$
    \item $\mathbb{R}^{n \times n}$ -- alle $n \times n$-Matrizen
    \item $\mathbb{R}[x]$ -- alle Polynome mit reellen Koeffizienten (z.B.\ $3x^2 + x - 5$)
    \item $\mathrm{FUNC}(\mathbb{R})$ -- alle Funktionen $f \colon \mathbb{R} \to \mathbb{R}$
\end{itemize}

Ein Vektorraum hat viele Regeln (Axiome): Assoziativit\"at, Kommutativit\"at, Nullelement, inverse Elemente, Distributivgesetze usw.\ -- insgesamt 8+ St\"uck. Die alle zu pr\"ufen w\"are m\"uhsam. Deshalb gibt es einen Trick...

%-----------------------------------------------------------------------------
\section*{2. Was ist ein Untervektorraum (Unterraum)?}

Ein \textbf{Untervektorraum} (kurz: Unterraum, abgek\"urzt UR) einer Menge $V$ ist eine \textbf{Teilmenge} $W \subseteq V$, die \textbf{selbst wieder ein Vektorraum} ist.

\begin{analogie}
Denk an eine \textbf{Firma} (= der gro\ss e Vektorraum $V$). Eine \textbf{Abteilung} innerhalb der Firma (= der Unterraum $W$) ist ein Teil der Firma. Aber die Abteilung muss alle Firmenregeln einhalten -- sie darf keine eigenen Regeln machen, die gegen die Firmenregeln versto\ss en.

Also: $W$ ist eine Teilmenge von $V$, und $W$ muss sich an alle Regeln von $V$ halten.
\end{analogie}

%-----------------------------------------------------------------------------
\section*{3. Der Unterraum-Test -- die 3 magischen Bedingungen}

\begin{merkbox}[title={Der Unterraum-Test (URT)}]
Sei $V$ ein Vektorraum und $W \subseteq V$. Dann ist $W$ ein Unterraum von $V$ genau dann, wenn:

\medskip
\textbf{(U1) Nullvektor:} $\vec{0} \in W$ \quad (``Der leere/neutrale Zustand muss dabei sein'')

\medskip
\textbf{(U2) Abgeschlossen unter Addition:}\\
F\"ur alle $u, v \in W$ gilt: $u + v \in W$\\
(``Wenn ich zwei Dinge aus $W$ addiere, muss das Ergebnis wieder in $W$ sein'')

\medskip
\textbf{(U3) Abgeschlossen unter Skalarmultiplikation:}\\
F\"ur alle $u \in W$ und alle $\lambda \in \mathbb{R}$ gilt: $\lambda \cdot u \in W$\\
(``Wenn ich etwas aus $W$ mit irgendeiner Zahl multipliziere, muss das Ergebnis wieder in $W$ sein'')
\end{merkbox}

\textbf{Warum reichen diese 3 Bedingungen?} Normalerweise m\"usste man ALLE 8+ Axiome eines Vektorraums pr\"ufen. Aber da $W$ eine Teilmenge von $V$ ist und $V$ schon ein Vektorraum ist, ``erbt'' $W$ automatisch die meisten Regeln von $V$:
\begin{itemize}
    \item Assoziativit\"at? Gilt schon in $V$ f\"ur alle Elemente, also auch f\"ur die in $W$.
    \item Kommutativit\"at? Gilt schon in $V$, also auch in $W$.
    \item Distributivgesetze? Genauso.
\end{itemize}

Die einzigen Dinge, die man EXTRA pr\"ufen muss, sind:
\begin{itemize}
    \item Ist das Nullelement dabei? (U1)
    \item Bleibt man ``drinnen'', wenn man addiert? (U2)
    \item Bleibt man ``drinnen'', wenn man skaliert? (U3)
\end{itemize}

\begin{analogie}
Die Firma hat schon alle Gesetze. F\"ur eine Abteilung musst du nur pr\"ufen: (1) Hat die Abteilung einen Chef (Nullelement/neutrales Element)? (2) Wenn zwei Leute in der Abteilung zusammenarbeiten, bleibt das Ergebnis in der Abteilung? (3) Wenn jemand aus der Abteilung eine Aufgabe skaliert (vergr\"o\ss ert/verkleinert), bleibt das in der Abteilung?
\end{analogie}

%-----------------------------------------------------------------------------
\section*{4. Strategie}

\begin{merkbox}[title={Strategie-\"Ubersicht}]
\textbf{Wenn du zeigen willst, dass $W$ ein Unterraum IST:}\\
$\rightarrow$ Beweise ALLE 3 Bedingungen (U1), (U2), (U3). Jede einzelne.

\medskip
\textbf{Wenn du zeigen willst, dass $W$ KEIN Unterraum ist:}\\
$\rightarrow$ Finde EIN konkretes Gegenbeispiel, das EINE der 3 Bedingungen verletzt.\\
(Tipp: Am einfachsten ist oft (U1) -- pr\"ufe, ob der Nullvektor drin ist!)
\end{merkbox}

%-----------------------------------------------------------------------------
\section*{5. Typische Fallen -- Woran erkenne ich sofort KEIN Unterraum?}

\begin{merkbox}[title={Rote Flaggen -- KEIN Unterraum}]
\textbf{Falle 1: Inhomogene Gleichungen}\\
Wenn die Bedingung so aussieht wie $a + 2c = b - 3$ (also eine Konstante $\neq 0$ auf einer Seite steht), dann ist der \textbf{Nullvektor meistens nicht drin}. Denn: Setzt man alle Variablen $= 0$, wird die Gleichung zu $0 = -3$, was falsch ist.

\medskip
Vergleiche: $a + 2c = b - 3d$ ist \textbf{homogen} (alle Terme enthalten Variablen, keine ``nackte'' Konstante). Setzt man alles $= 0$: $0 = 0$ \checkmark

\bigskip
\textbf{Falle 2: Ungleichungen ($\leq$, $<$, $\geq$, $>$)}\\
Meistens \textbf{nicht abgeschlossen unter Skalarmultiplikation}. Warum? Wenn du mit $-1$ multiplizierst, dreht sich das Vorzeichen um, und $\leq$ wird zu $\geq$!

\bigskip
\textbf{Falle 3: Produkte, Quadrate, Betr\"age}\\
Bedingungen wie $ac \leq 0$, $A^2 = 0$, $|a| \leq |c|$, $p(1) \cdot p(2) = 0$ sind meistens \textbf{nicht abgeschlossen unter Addition}. Nichtlineare Bedingungen vertragen sich schlecht mit Addition.

\bigskip
\textbf{Falle 4: $\neq$ (Ungleich)}\\
Wenn die Bedingung ein $\neq$ enth\"alt, ist der Nullvektor oft nicht drin, oder die Abgeschlossenheit scheitert.
\end{merkbox}

\begin{merkbox}[title={Gr\"une Flagge -- JA, Unterraum!}]
\textbf{Homogene lineare Bedingungen sind IMMER Untervektorr\"aume.}\\
Das hei\ss t: Wenn die Bedingung eine lineare Gleichung ist, bei der auf einer Seite $0$ steht (oder sich alles zu $0$ vereinfachen l\"asst), dann ist es ein Unterraum.

\medskip
Beispiele: $a + 2c = b - 3d$ (\"aquivalent zu $a - b + 2c + 3d = 0$), $A^\top = -A$, $p(1) = 0$, $f(-x) = -f(x)$.
\end{merkbox}

\newpage

%=============================================================================
%  TEIL 2: AUFGABE 10 (R^n Unterräume)
%=============================================================================
\part*{Teil 2: Aufgabe 10 -- Untervektorr\"aume in $\mathbb{R}^n$}

Hier ist $V = \mathbb{R}^n$ (also Vektoren aus reellen Zahlen) und wir pr\"ufen verschiedene Teilmengen $W$.

%-----------------------------------------------------------------------------
\subsection*{(a) \quad $V = \mathbb{R}^4$, \quad $W = \{(a,b,c,d) \in \mathbb{R}^4 : a + 2c = b - 3d\}$}

\textbf{Was bedeutet die Bedingung?}\\
Wir nehmen alle Vektoren $(a,b,c,d)$ aus $\mathbb{R}^4$, bei denen die Gleichung $a + 2c = b - 3d$ erf\"ullt ist. Umgeschrieben: $a - b + 2c + 3d = 0$. Das ist eine \textbf{homogene lineare Gleichung} (rechts steht $0$, links stehen nur Variablen-Terme).

\begin{merkbox}
Homogene lineare Gleichung $\Rightarrow$ immer Unterraum! Aber wir beweisen es trotzdem sauber.
\end{merkbox}

\textbf{Beweis: $W$ ist ein Unterraum.}

\textbf{(U1) Nullvektor:} Setze $a = b = c = d = 0$: $0 + 0 = 0 - 0 \Leftrightarrow 0 = 0$ \checkmark\\
Also $(0,0,0,0) \in W$.

\textbf{(U2) Addition:} Seien $u = (a_1, b_1, c_1, d_1) \in W$ und $v = (a_2, b_2, c_2, d_2) \in W$.\\
Das hei\ss t: $a_1 + 2c_1 = b_1 - 3d_1$ und $a_2 + 2c_2 = b_2 - 3d_2$.

$u + v = (a_1 + a_2,\; b_1 + b_2,\; c_1 + c_2,\; d_1 + d_2)$.

Pr\"ufe die Bedingung f\"ur $u + v$:
\begin{align*}
(a_1 + a_2) + 2(c_1 + c_2) &= (a_1 + 2c_1) + (a_2 + 2c_2) \\
&= (b_1 - 3d_1) + (b_2 - 3d_2) \quad \text{(weil $u, v \in W$)} \\
&= (b_1 + b_2) - 3(d_1 + d_2) \quad \checkmark
\end{align*}
Also $u + v \in W$.

\textbf{(U3) Skalarmultiplikation:} Sei $u = (a_1, b_1, c_1, d_1) \in W$ und $\lambda \in \mathbb{R}$.\\
Also $a_1 + 2c_1 = b_1 - 3d_1$.

$\lambda u = (\lambda a_1, \lambda b_1, \lambda c_1, \lambda d_1)$.

Pr\"ufe:
\begin{align*}
\lambda a_1 + 2 \lambda c_1 &= \lambda(a_1 + 2c_1) = \lambda(b_1 - 3d_1) = \lambda b_1 - 3\lambda d_1 \quad \checkmark
\end{align*}
Also $\lambda u \in W$.

\medskip
$\Rightarrow$ \textbf{$W$ ist ein Unterraum von $\mathbb{R}^4$.} \hfill $\blacksquare$

%-----------------------------------------------------------------------------
\subsection*{(b) \quad $V = \mathbb{R}^4$, \quad $W = \{(a,b,c,d) \in \mathbb{R}^4 : a + 2c = b - 3\}$}

\textbf{Was bedeutet die Bedingung?}\\
Fast wie (a), ABER rechts steht $b - 3$ statt $b - 3d$. Die $-3$ ist eine \textbf{Konstante}, nicht $-3d$ (ein Variablen-Term). Umgeschrieben: $a - b + 2c = -3$. Das ist eine \textbf{inhomogene} Gleichung (rechts steht $-3 \neq 0$).

\begin{merkbox}
Inhomogene Gleichung $\Rightarrow$ Nullvektor wahrscheinlich nicht drin $\Rightarrow$ kein Unterraum!
\end{merkbox}

\textbf{$W$ ist KEIN Unterraum.}

\textbf{Gegenbeispiel (U1 verletzt):} Teste den Nullvektor $(0,0,0,0)$:
\[
0 + 2 \cdot 0 = 0 - 3 \quad \Leftrightarrow \quad 0 = -3 \quad \text{\Lightning}
\]
Der Nullvektor erf\"ullt die Bedingung \textbf{nicht}. Also $(0,0,0,0) \notin W$.

$\Rightarrow$ \textbf{$W$ ist kein Unterraum.} \hfill $\blacksquare$

%-----------------------------------------------------------------------------
\subsection*{(c) \quad $V = \mathbb{R}^3$, \quad $W = \{(a,b,c) \in \mathbb{R}^3 : a \leq c\}$}

\textbf{Was bedeutet die Bedingung?}\\
Alle Vektoren $(a,b,c)$, bei denen die erste Komponente $a$ kleiner oder gleich der dritten Komponente $c$ ist. Zum Beispiel: $(1, 5, 3) \in W$ weil $1 \leq 3$. Aber $(4, 5, 2) \notin W$ weil $4 \not\leq 2$.

\begin{merkbox}
Ungleichung ($\leq$) $\Rightarrow$ Vorsicht bei Skalarmultiplikation mit negativen Zahlen!
\end{merkbox}

\textbf{$W$ ist KEIN Unterraum.}

\textbf{Gegenbeispiel (U3 verletzt):} Nimm $u = (1, 0, 2) \in W$ (denn $1 \leq 2$ \checkmark) und $\lambda = -1$.

\[
\lambda \cdot u = -1 \cdot (1, 0, 2) = (-1, 0, -2)
\]

Pr\"ufe: $a \leq c$? $\quad -1 \leq -2$? $\quad$ \textbf{NEIN!} ($-1 > -2$) \quad \Lightning

\textbf{Warum passiert das?} Wenn man mit $-1$ multipliziert, drehen sich die Vorzeichen um. Aus $1 \leq 2$ wird $-1 \leq -2$, was falsch ist. Die $\leq$-Richtung ``kippt'' bei negativer Skalarmultiplikation!

$\Rightarrow$ \textbf{$W$ ist kein Unterraum.} \hfill $\blacksquare$

%-----------------------------------------------------------------------------
\subsection*{(d) \quad $V = \mathbb{R}^3$, \quad $W = \{(a,b,c) \in \mathbb{R}^3 : |a| \leq |c|\}$}

\textbf{Was bedeutet die Bedingung?}\\
Der \textbf{Betrag} der ersten Komponente muss kleiner oder gleich dem Betrag der dritten sein. Zum Beispiel: $(-2, 0, 3) \in W$ weil $|-2| = 2 \leq 3 = |3|$.

Betr\"age machen die Sache nichtlinear -- man k\"onnte denken, dass der Betrag das Problem mit negativer Skalarmultiplikation l\"ost (weil $|-a| = |a|$). Aber das Problem verschiebt sich zur Addition!

\textbf{$W$ ist KEIN Unterraum.}

\textbf{Gegenbeispiel (U2 verletzt):} Nimm:
\[
u = (1, 0, -1) \in W \quad \text{(denn } |1| = 1 \leq 1 = |-1| \text{ \checkmark)}
\]
\[
v = (1, 0, 1) \in W \quad \text{(denn } |1| = 1 \leq 1 = |1| \text{ \checkmark)}
\]
\[
u + v = (2, 0, 0)
\]

Pr\"ufe: $|2| \leq |0|$? $\quad 2 \leq 0$? $\quad$ \textbf{NEIN!} \quad \Lightning

\textbf{Warum passiert das?} Die Betr\"age der dritten Komponenten haben sich \textbf{weggek\"urzt} ($-1 + 1 = 0$), w\"ahrend die Betr\"age der ersten sich \textbf{addiert} haben ($1 + 1 = 2$). Betr\"age und Addition vertragen sich nicht gut.

$\Rightarrow$ \textbf{$W$ ist kein Unterraum.} \hfill $\blacksquare$

%-----------------------------------------------------------------------------
\subsection*{(e) \quad $V = \mathbb{R}^3$, \quad $W = \{(a,b,c) \in \mathbb{R}^3 : ac \leq 0\}$}

\textbf{Was bedeutet die Bedingung?}\\
Das \textbf{Produkt} der ersten und dritten Komponente muss $\leq 0$ sein. Das hei\ss t: $a$ und $c$ m\"ussen \textbf{verschiedene Vorzeichen} haben (oder mindestens einer ist $0$). Zum Beispiel:
\begin{itemize}
    \item $(1, 0, -2) \in W$ weil $1 \cdot (-2) = -2 \leq 0$ \checkmark
    \item $(3, 0, 5) \notin W$ weil $3 \cdot 5 = 15 > 0$ \Lightning
\end{itemize}

\begin{merkbox}
Produkt-Bedingungen ($ac \leq 0$) sind nichtlinear $\Rightarrow$ Addition geht meistens schief!
\end{merkbox}

\textbf{$W$ ist KEIN Unterraum.}

\textbf{Gegenbeispiel (U2 verletzt):} Nimm:
\[
u = (1, 0, 0) \in W \quad \text{(denn } 1 \cdot 0 = 0 \leq 0 \text{ \checkmark)}
\]
\[
v = (0, 0, 1) \in W \quad \text{(denn } 0 \cdot 1 = 0 \leq 0 \text{ \checkmark)}
\]
\[
u + v = (1, 0, 1)
\]

Pr\"ufe: $ac = 1 \cdot 1 = 1 \leq 0$? $\quad$ \textbf{NEIN!} ($1 > 0$) \quad \Lightning

\textbf{Warum passiert das?} Einzeln waren $a$ und $c$ ``brav'' (eines war immer $0$). Aber nach der Addition sind pl\"otzlich beide positiv, und das Produkt wird positiv.

$\Rightarrow$ \textbf{$W$ ist kein Unterraum.} \hfill $\blacksquare$

%-----------------------------------------------------------------------------
\subsection*{(f) \quad $V = \mathbb{R}^3$, \quad $W = \{(a,b,c) \in \mathbb{R}^3 : a + b \neq c + 1\}$}

\textbf{Was bedeutet die Bedingung?}\\
Alle Vektoren, bei denen $a + b$ \textbf{nicht gleich} $c + 1$ ist. Also: fast alle Vektoren AUSSER die, wo $a + b = c + 1$ gilt.

\textbf{$W$ ist KEIN Unterraum.}

\textbf{Gegenbeispiel (U2 verletzt):} Nimm:
\[
u = (2, 0, 0) \in W \quad \text{(denn } 2 + 0 = 2 \neq 0 + 1 = 1 \text{ \checkmark)}
\]
\[
v = (0, 0, 1) \in W \quad \text{(denn } 0 + 0 = 0 \neq 1 + 1 = 2 \text{ \checkmark)}
\]
\[
u + v = (2, 0, 1)
\]

Pr\"ufe: $a + b \neq c + 1$? $\quad 2 + 0 \neq 1 + 1$? $\quad 2 \neq 2$? $\quad$ \textbf{DOCH, 2 = 2!} \quad \Lightning

Also $u + v \notin W$.

\textbf{Warum passiert das?} Man kann immer zwei Vektoren so w\"ahlen, dass sie sich zu einem addieren, der genau die ``verbotene'' Gleichung erf\"ullt.

\textbf{Hinweis:} Man h\"atte auch (U1) pr\"ufen k\"onnen: $(0,0,0)$ -- ist $0 + 0 \neq 0 + 1$? $0 \neq 1$? Ja, also $(0,0,0) \in W$, (U1) ist erf\"ullt. Deshalb brauchten wir hier (U2).

$\Rightarrow$ \textbf{$W$ ist kein Unterraum.} \hfill $\blacksquare$

%-----------------------------------------------------------------------------
\subsection*{(g) \quad $V = \mathbb{R}^4$, \quad $W = \{\ell \in \mathbb{R}^4 : \ell \text{ ist eine arithmetische Folge}\}$}

\textbf{Erstmal: Was ist eine arithmetische Folge?}\\
Eine Folge $(a, b, c, d)$ hei\ss t \textbf{arithmetisch}, wenn der \textbf{Abstand} zwischen aufeinanderfolgenden Gliedern immer gleich ist.

Beispiele:
\begin{itemize}
    \item $(2, 5, 8, 11)$ -- Abstand immer $+3$ \checkmark
    \item $(10, 7, 4, 1)$ -- Abstand immer $-3$ \checkmark
    \item $(0, 0, 0, 0)$ -- Abstand immer $0$ \checkmark
    \item $(1, 2, 4, 8)$ -- Abst\"ande: $1, 2, 4$ -- NICHT gleich \Lightning
\end{itemize}

\textbf{Mathematisch:} $(a, b, c, d)$ ist arithmetisch $\Leftrightarrow$ $b - a = c - b = d - c$.

Das kann man umschreiben zu zwei Gleichungen:
\[
b - a = c - b \quad \Leftrightarrow \quad a - 2b + c = 0
\]
\[
c - b = d - c \quad \Leftrightarrow \quad b - 2c + d = 0
\]

\begin{merkbox}
Das sind zwei \textbf{homogene lineare Gleichungen}! (Rechts steht $0$.) Also ist $W$ ein Unterraum!
\end{merkbox}

\textbf{Beweis: $W$ ist ein Unterraum.}

\textbf{(U1) Nullvektor:} $(0, 0, 0, 0)$: $0 - 2 \cdot 0 + 0 = 0$ \checkmark{} und $0 - 2 \cdot 0 + 0 = 0$ \checkmark. Also $(0,0,0,0) \in W$.

\textbf{(U2) Addition:} Seien $u = (a_1, b_1, c_1, d_1)$ und $v = (a_2, b_2, c_2, d_2)$ in $W$.\\
Dann: $a_1 - 2b_1 + c_1 = 0$, $b_1 - 2c_1 + d_1 = 0$, $a_2 - 2b_2 + c_2 = 0$, $b_2 - 2c_2 + d_2 = 0$.

F\"ur $u + v = (a_1+a_2, b_1+b_2, c_1+c_2, d_1+d_2)$:
\begin{align*}
(a_1+a_2) - 2(b_1+b_2) + (c_1+c_2) &= (a_1 - 2b_1 + c_1) + (a_2 - 2b_2 + c_2) = 0 + 0 = 0 \quad \checkmark \\
(b_1+b_2) - 2(c_1+c_2) + (d_1+d_2) &= (b_1 - 2c_1 + d_1) + (b_2 - 2c_2 + d_2) = 0 + 0 = 0 \quad \checkmark
\end{align*}

\textbf{(U3) Skalarmultiplikation:} Sei $u = (a, b, c, d) \in W$, $\lambda \in \mathbb{R}$.
\[
\lambda a - 2\lambda b + \lambda c = \lambda(a - 2b + c) = \lambda \cdot 0 = 0 \quad \checkmark
\]
Genauso f\"ur die zweite Gleichung.

$\Rightarrow$ \textbf{$W$ ist ein Unterraum von $\mathbb{R}^4$.} \hfill $\blacksquare$

%-----------------------------------------------------------------------------
\subsection*{(h) \quad $V = \mathbb{R}^4$, \quad $W = \{\ell \in \mathbb{R}^4 : \ell \text{ ist eine geometrische Folge}\}$}

\textbf{Was ist eine geometrische Folge?}\\
Eine Folge $(a, b, c, d)$ hei\ss t \textbf{geometrisch}, wenn das \textbf{Verh\"altnis} zwischen aufeinanderfolgenden Gliedern immer gleich ist.

Beispiele:
\begin{itemize}
    \item $(1, 2, 4, 8)$ -- Verh\"altnis immer $\times 2$ \checkmark
    \item $(1, 3, 9, 27)$ -- Verh\"altnis immer $\times 3$ \checkmark
    \item $(81, 27, 9, 3)$ -- Verh\"altnis immer $\times \frac{1}{3}$ \checkmark
\end{itemize}

\textbf{Mathematisch:} $\frac{b}{a} = \frac{c}{b} = \frac{d}{c}$ (also $b^2 = ac$ und $c^2 = bd$). Das sind \textbf{quadratische} (nichtlineare!) Bedingungen.

\textbf{$W$ ist KEIN Unterraum.}

\textbf{Gegenbeispiel (U2 verletzt):}
\[
u = (1, 2, 4, 8) \in W \quad \text{(Verh\"altnis $= 2$)}
\]
\[
v = (1, 3, 9, 27) \in W \quad \text{(Verh\"altnis $= 3$)}
\]
\[
u + v = (2, 5, 13, 35)
\]

Ist das eine geometrische Folge? Pr\"ufe: $\frac{5}{2} = 2{,}5$ aber $\frac{13}{5} = 2{,}6$. Die Verh\"altnisse sind \textbf{nicht gleich}! \Lightning

\textbf{Warum passiert das?} Geometrische Folgen sind durch \textbf{Multiplikation} definiert (Verh\"altnisse), aber wir \textbf{addieren} die Vektoren. Multiplikation und Addition vertragen sich hier nicht.

$\Rightarrow$ \textbf{$W$ ist kein Unterraum.} \hfill $\blacksquare$

%-----------------------------------------------------------------------------
\subsection*{(i) \quad $V = \mathbb{R}^3$, \quad $W = \mathbb{R}^2$}

\textbf{Was bedeutet das?}\\
Hier ist $V = \mathbb{R}^3$ (Vektoren mit 3 Komponenten) und $W = \mathbb{R}^2$ (Vektoren mit 2 Komponenten).

\textbf{$W$ ist KEIN Unterraum.}

\textbf{Warum?} Der allererste Schritt: $W$ muss eine \textbf{Teilmenge} von $V$ sein ($W \subseteq V$). Aber:
\begin{itemize}
    \item Ein Element von $\mathbb{R}^2$ sieht so aus: $(x, y)$ -- zwei Komponenten.
    \item Ein Element von $\mathbb{R}^3$ sieht so aus: $(x, y, z)$ -- drei Komponenten.
\end{itemize}

Das sind \textbf{verschiedene Typen von Objekten}! $(1, 2) \in \mathbb{R}^2$, aber $(1, 2) \notin \mathbb{R}^3$, weil in $\mathbb{R}^3$ jeder Vektor \textbf{drei} Komponenten braucht.

\begin{analogie}
Das ist so, als w\"urdest du sagen: ``Die Abteilung Fische ist ein Teil der Firma Autos.'' Das geht nicht, weil Fische keine Autos sind -- sie sind etwas komplett anderes.
\end{analogie}

Also: $\mathbb{R}^2 \not\subseteq \mathbb{R}^3$, und damit kann $W$ kein Unterraum sein.

$\Rightarrow$ \textbf{$W$ ist kein Unterraum von $V$.} \hfill $\blacksquare$

\newpage

%=============================================================================
%  TEIL 3: AUFGABE 11 (Matrizen-Unterräume)
%=============================================================================
\part*{Teil 3: Aufgabe 11 -- Untervektorr\"aume in $\mathbb{R}^{n \times n}$}

Hier ist $V = \mathbb{R}^{n \times n}$ mit $n > 1$, also der Vektorraum aller $n \times n$-Matrizen. Die Vektoren sind jetzt \textbf{Matrizen}! Der Nullvektor ist die \textbf{Nullmatrix} $0$ (alle Eintr\"age $= 0$).

\textbf{Kurze Erinnerung an Begriffe:}

\begin{merkbox}[title={Wichtige Matrixbegriffe}]
\textbf{Transponierte $A^\top$:} Zeilen und Spalten tauschen. Also: $(A^\top)_{ij} = A_{ji}$.\\
Beispiel: $\begin{pmatrix} 1 & 2 \\ 3 & 4 \end{pmatrix}^\top = \begin{pmatrix} 1 & 3 \\ 2 & 4 \end{pmatrix}$

\medskip
\textbf{Schiefsymmetrisch:} $A^\top = -A$. Die Diagonale ist immer $0$, und die Eintr\"age ``spiegeln sich mit Vorzeichenwechsel''.\\
Beispiel: $\begin{pmatrix} 0 & 3 \\ -3 & 0 \end{pmatrix}$ -- oben rechts $3$, unten links $-3$.

\medskip
\textbf{Spur (Trace):} Die Summe der Diagonaleintr\"age: $\mathrm{Spur}(A) = a_{11} + a_{22} + \ldots + a_{nn}$.\\
Beispiel: $\mathrm{Spur}\begin{pmatrix} 1 & 2 \\ 3 & 4 \end{pmatrix} = 1 + 4 = 5$.

\medskip
\textbf{Stufenmatrix (Zeilenstufenform):} Jede Zeile beginnt weiter rechts als die dar\"uber, und Nullzeilen stehen ganz unten.\\
Beispiel: $\begin{pmatrix} 1 & 2 & 3 \\ 0 & 1 & 5 \\ 0 & 0 & 1 \end{pmatrix}$ -- jede ``Stufe'' geht nach rechts.
\end{merkbox}

%-----------------------------------------------------------------------------
\subsection*{(a) \quad $W$ = Menge aller schiefsymmetrischen Matrizen ($A^\top = -A$)}

\textbf{Was bedeutet die Bedingung?}\\
$A^\top = -A$ hei\ss t: Wenn man die Matrix transponiert (Zeilen $\leftrightarrow$ Spalten), bekommt man $-A$. Die Diagonale muss $0$ sein (denn $a_{ii} = -a_{ii}$ geht nur f\"ur $a_{ii} = 0$), und die Eintr\"age au\ss erhalb der Diagonalen spiegeln sich mit Vorzeichenwechsel.

\begin{merkbox}
$A^\top = -A$ ist eine \textbf{lineare Bedingung} an die Eintr\"age von $A$. Also erwarten wir: JA, Unterraum!
\end{merkbox}

\textbf{Beweis: $W$ ist ein Unterraum.}

\textbf{(U1) Nullmatrix:} $0^\top = 0 = -0$ \checkmark. Also $0 \in W$.

\textbf{(U2) Addition:} Seien $A, B \in W$, also $A^\top = -A$ und $B^\top = -B$.
\[
(A + B)^\top = A^\top + B^\top = (-A) + (-B) = -(A + B) \quad \checkmark
\]
(Wir nutzen hier: die Transponierte einer Summe ist die Summe der Transponierten.)

\textbf{(U3) Skalarmultiplikation:} Sei $A \in W$, $\lambda \in \mathbb{R}$.
\[
(\lambda A)^\top = \lambda \cdot A^\top = \lambda \cdot (-A) = -(\lambda A) \quad \checkmark
\]

$\Rightarrow$ \textbf{$W$ ist ein Unterraum von $\mathbb{R}^{n \times n}$.} \hfill $\blacksquare$

%-----------------------------------------------------------------------------
\subsection*{(b) \quad $W = \{A \in \mathbb{R}^{n \times n} : A^2 = 0\}$}

\textbf{Was bedeutet die Bedingung?}\\
Wenn man $A$ mit sich selbst multipliziert ($A \cdot A$), kommt die Nullmatrix raus. Solche Matrizen hei\ss en \textbf{nilpotent} (mit Nilpotenzgrad $\leq 2$).

\begin{merkbox}
$A^2 = 0$ ist eine \textbf{quadratische} (nichtlineare!) Bedingung. Also erwarten wir Probleme bei der Addition!
\end{merkbox}

\textbf{$W$ ist KEIN Unterraum.}

\textbf{Gegenbeispiel (U2 verletzt) f\"ur $n = 2$:} Nimm:
\[
A = \begin{pmatrix} 0 & 1 \\ 0 & 0 \end{pmatrix}, \quad
B = \begin{pmatrix} 0 & 0 \\ 1 & 0 \end{pmatrix}
\]

Pr\"ufe: $A \in W$?
\[
A^2 = \begin{pmatrix} 0 & 1 \\ 0 & 0 \end{pmatrix} \begin{pmatrix} 0 & 1 \\ 0 & 0 \end{pmatrix} = \begin{pmatrix} 0 \cdot 0 + 1 \cdot 0 & 0 \cdot 1 + 1 \cdot 0 \\ 0 \cdot 0 + 0 \cdot 0 & 0 \cdot 1 + 0 \cdot 0 \end{pmatrix} = \begin{pmatrix} 0 & 0 \\ 0 & 0 \end{pmatrix} \quad \checkmark
\]

Pr\"ufe: $B \in W$?
\[
B^2 = \begin{pmatrix} 0 & 0 \\ 1 & 0 \end{pmatrix} \begin{pmatrix} 0 & 0 \\ 1 & 0 \end{pmatrix} = \begin{pmatrix} 0 & 0 \\ 0 & 0 \end{pmatrix} \quad \checkmark
\]

Jetzt die Summe:
\[
A + B = \begin{pmatrix} 0 & 1 \\ 1 & 0 \end{pmatrix}
\]

\[
(A + B)^2 = \begin{pmatrix} 0 & 1 \\ 1 & 0 \end{pmatrix} \begin{pmatrix} 0 & 1 \\ 1 & 0 \end{pmatrix} = \begin{pmatrix} 0 \cdot 0 + 1 \cdot 1 & 0 \cdot 1 + 1 \cdot 0 \\ 1 \cdot 0 + 0 \cdot 1 & 1 \cdot 1 + 0 \cdot 0 \end{pmatrix} = \begin{pmatrix} 1 & 0 \\ 0 & 1 \end{pmatrix} = I \neq 0 \quad \Lightning
\]

\textbf{Warum passiert das?} $(A+B)^2 = A^2 + AB + BA + B^2$. Wir wissen $A^2 = 0$ und $B^2 = 0$, also $(A+B)^2 = AB + BA$. Aber $AB + BA$ muss NICHT $= 0$ sein! Bei Matrizen ist $AB \neq BA$ im Allgemeinen, also k\"urzt sich nichts weg.

$\Rightarrow$ \textbf{$W$ ist kein Unterraum.} \hfill $\blacksquare$

%-----------------------------------------------------------------------------
\subsection*{(c) \quad $W$ = Menge aller $A \in \mathbb{R}^{n \times n}$ mit $A \begin{pmatrix} 1 \\ 2 \\ \vdots \\ n \end{pmatrix} = \begin{pmatrix} 0 \\ 0 \\ \vdots \\ 0 \end{pmatrix}$}

\textbf{Was bedeutet die Bedingung?}\\
Wir nennen $w = (1, 2, \ldots, n)^\top$. Die Bedingung ist: $Aw = 0$. Das hei\ss t: wenn man $A$ auf den speziellen Vektor $w$ anwendet, kommt der Nullvektor raus. $W$ ist also der \textbf{Kern} der Abbildung ``multipliziere mit $w$''.

\begin{merkbox}
$Aw = 0$ ist eine \textbf{homogene lineare Bedingung} an $A$ (man multipliziert $A$ mit einem festen Vektor und fordert, dass $0$ rauskommt). Also: JA, Unterraum!
\end{merkbox}

\textbf{Beweis: $W$ ist ein Unterraum.}

\textbf{(U1)} $0 \cdot w = 0$ \checkmark. Die Nullmatrix multipliziert mit allem ergibt $0$.

\textbf{(U2)} Seien $A, B \in W$, also $Aw = 0$ und $Bw = 0$.
\[
(A + B)w = Aw + Bw = 0 + 0 = 0 \quad \checkmark
\]

\textbf{(U3)} Sei $A \in W$, $\lambda \in \mathbb{R}$.
\[
(\lambda A)w = \lambda (Aw) = \lambda \cdot 0 = 0 \quad \checkmark
\]

$\Rightarrow$ \textbf{$W$ ist ein Unterraum von $\mathbb{R}^{n \times n}$.} \hfill $\blacksquare$

%-----------------------------------------------------------------------------
\subsection*{(d) \quad $W$ = Menge aller Stufenmatrizen in $\mathbb{R}^{n \times n}$}

\textbf{Was bedeutet die Bedingung?}\\
$W$ enth\"alt alle Matrizen, die in \textbf{Zeilenstufenform} sind. In der Zeilenstufenform beginnt jede nicht-Null-Zeile weiter rechts als die dar\"uber, und Nullzeilen stehen unten.

\textbf{$W$ ist KEIN Unterraum.}

\textbf{Gegenbeispiel (U2 verletzt) f\"ur $n = 2$:}

\[
A = \begin{pmatrix} 1 & 0 \\ 0 & 1 \end{pmatrix} \quad \text{(Stufenform \checkmark -- das ist sogar die Einheitsmatrix)}
\]
\[
B = \begin{pmatrix} -1 & 0 \\ 0 & 0 \end{pmatrix} \quad \text{(Stufenform \checkmark -- zweite Zeile ist Nullzeile)}
\]
\[
A + B = \begin{pmatrix} 0 & 0 \\ 0 & 1 \end{pmatrix}
\]

Ist $A + B$ in Stufenform? Die erste Zeile ist $(0, 0)$ (Nullzeile), die zweite Zeile ist $(0, 1)$ (nicht-Nullzeile). Aber in Stufenform m\"ussen Nullzeilen \textbf{unten} stehen! Hier steht die Nullzeile \textbf{oben}. \Lightning

\textbf{Warum passiert das?} Stufenform h\"angt von der \textbf{Position} der f\"uhrenden Eintr\"age ab. Addition kann die f\"uhrenden Eintr\"age wegl\"oschen und damit die Stufenstruktur zerst\"oren.

$\Rightarrow$ \textbf{$W$ ist kein Unterraum.} \hfill $\blacksquare$

%-----------------------------------------------------------------------------
\subsection*{(e) \quad $W = \{A \in \mathbb{R}^{n \times n} : KAM = A^\top\}$, wobei $K, M$ fest sind}

\textbf{Was bedeutet die Bedingung?}\\
$K$ und $M$ sind zwei feste (gegebene) $n \times n$-Matrizen. Die Bedingung sagt: Wenn man $A$ von links mit $K$ und von rechts mit $M$ multipliziert, kommt die Transponierte $A^\top$ raus.

Umgeschrieben: $KAM - A^\top = 0$. Das ist eine \textbf{lineare Bedingung} an $A$ (weil Matrixmultiplikation und Transponierung linear in $A$ sind).

\begin{merkbox}
Lineare Bedingung $= 0$ $\Rightarrow$ Unterraum!
\end{merkbox}

\textbf{Beweis: $W$ ist ein Unterraum.}

\textbf{(U1) Nullmatrix:} $K \cdot 0 \cdot M = 0 = 0^\top$ \checkmark.

\textbf{(U2) Addition:} Seien $A, B \in W$, also $KAM = A^\top$ und $KBM = B^\top$.
\[
K(A+B)M = KAM + KBM = A^\top + B^\top = (A+B)^\top \quad \checkmark
\]

\textbf{(U3) Skalarmultiplikation:} Sei $A \in W$, $\lambda \in \mathbb{R}$.
\[
K(\lambda A)M = \lambda (KAM) = \lambda A^\top = (\lambda A)^\top \quad \checkmark
\]

$\Rightarrow$ \textbf{$W$ ist ein Unterraum von $\mathbb{R}^{n \times n}$.} \hfill $\blacksquare$

%-----------------------------------------------------------------------------
\subsection*{(f) \quad $W$ = Menge aller $A \in \mathbb{R}^{n \times n}$ mit Spur$(A) = 0$}

\textbf{Was bedeutet die Bedingung?}\\
Die \textbf{Spur} einer Matrix ist die Summe der Diagonaleintr\"age. Wir wollen alle Matrizen, deren Diagonalsumme $0$ ist.

Zum Beispiel f\"ur $n = 2$: $A = \begin{pmatrix} 3 & 7 \\ -1 & -3 \end{pmatrix}$ hat $\mathrm{Spur}(A) = 3 + (-3) = 0$ \checkmark.

\begin{merkbox}
Die Spur ist eine \textbf{lineare Funktion}: $\mathrm{Spur}(A + B) = \mathrm{Spur}(A) + \mathrm{Spur}(B)$ und $\mathrm{Spur}(\lambda A) = \lambda \cdot \mathrm{Spur}(A)$. Lineare Bedingung $= 0$ $\Rightarrow$ Unterraum!
\end{merkbox}

\textbf{Beweis: $W$ ist ein Unterraum.}

\textbf{(U1) Nullmatrix:} $\mathrm{Spur}(0) = 0 + 0 + \ldots + 0 = 0$ \checkmark.

\textbf{(U2) Addition:} Seien $A, B \in W$, also $\mathrm{Spur}(A) = 0$ und $\mathrm{Spur}(B) = 0$.
\[
\mathrm{Spur}(A + B) = \mathrm{Spur}(A) + \mathrm{Spur}(B) = 0 + 0 = 0 \quad \checkmark
\]

\textbf{(U3) Skalarmultiplikation:} Sei $A \in W$, $\lambda \in \mathbb{R}$.
\[
\mathrm{Spur}(\lambda A) = \lambda \cdot \mathrm{Spur}(A) = \lambda \cdot 0 = 0 \quad \checkmark
\]

$\Rightarrow$ \textbf{$W$ ist ein Unterraum von $\mathbb{R}^{n \times n}$.} \hfill $\blacksquare$

\newpage

%=============================================================================
%  TEIL 4: AUFGABE 12 (Polynom-Unterräume)
%=============================================================================
\part*{Teil 4: Aufgabe 12 -- Untervektorr\"aume in $\mathbb{R}[x]$ (Polynome)}

Hier ist $V = \mathbb{R}[x]$, der Vektorraum aller Polynome mit reellen Koeffizienten. Die ``Vektoren'' sind jetzt Polynome wie $3x^2 + x - 5$ oder $x^{100}$ oder einfach die Zahl $7$. Der Nullvektor ist das \textbf{Nullpolynom} $p(x) = 0$ (das Polynom, das \"uberall $0$ ist).

\textbf{Wichtig:} Polynome kann man addieren und mit Zahlen multiplizieren -- genau wie normale Vektoren!

\textbf{Kurze Erinnerung:}

\begin{merkbox}[title={Polynome -- Basics}]
\textbf{Was ist $p(1)$?} Setze $x = 1$ ins Polynom ein. Z.B.\ $p(x) = x^2 + 3x$, dann $p(1) = 1 + 3 = 4$.

\medskip
\textbf{Was ist $p'(x)$?} Die \textbf{Ableitung} von $p(x)$. Z.B.\ $p(x) = x^3 + 2x$, dann $p'(x) = 3x^2 + 2$.\\
Regeln: $(x^n)' = n \cdot x^{n-1}$, (Konstante)$' = 0$, Ableitung ist linear.
\end{merkbox}

%-----------------------------------------------------------------------------
\subsection*{(a) \quad $W = \{p(x) \in \mathbb{R}[x] : p(1) = 0\}$}

\textbf{Was bedeutet die Bedingung?}\\
Alle Polynome, die an der Stelle $x = 1$ den Wert $0$ haben. Also: $x = 1$ ist eine \textbf{Nullstelle} des Polynoms.

Beispiele:
\begin{itemize}
    \item $p(x) = x - 1 \in W$ \quad (denn $p(1) = 0$)
    \item $p(x) = x^2 - 1 \in W$ \quad (denn $p(1) = 0$)
    \item $p(x) = x \notin W$ \quad (denn $p(1) = 1 \neq 0$)
\end{itemize}

\begin{merkbox}
Die Abbildung $p \mapsto p(1)$ ist \textbf{linear}: $(p + q)(1) = p(1) + q(1)$ und $(\lambda p)(1) = \lambda \cdot p(1)$.\\
Lineare Bedingung $= 0$ $\Rightarrow$ Unterraum!
\end{merkbox}

\textbf{Beweis: $W$ ist ein Unterraum.}

\textbf{(U1) Nullpolynom:} Das Nullpolynom $p(x) = 0$ hat $p(1) = 0$ \checkmark.

\textbf{(U2) Addition:} Seien $p, q \in W$, also $p(1) = 0$ und $q(1) = 0$.
\[
(p + q)(1) = p(1) + q(1) = 0 + 0 = 0 \quad \checkmark
\]

\textbf{(U3) Skalarmultiplikation:} Sei $p \in W$, $\lambda \in \mathbb{R}$.
\[
(\lambda p)(1) = \lambda \cdot p(1) = \lambda \cdot 0 = 0 \quad \checkmark
\]

$\Rightarrow$ \textbf{$W$ ist ein Unterraum von $\mathbb{R}[x]$.} \hfill $\blacksquare$

%-----------------------------------------------------------------------------
\subsection*{(b) \quad $W = \{p(x) \in \mathbb{R}[x] : p(1) = p(2) = 0\}$}

\textbf{Was bedeutet die Bedingung?}\\
Alle Polynome, die SOWOHL bei $x = 1$ als auch bei $x = 2$ den Wert $0$ haben. Also: $1$ und $2$ sind beides Nullstellen.

Das ist einfach die \textbf{Kombination} von zwei linearen Bedingungen: $p(1) = 0$ UND $p(2) = 0$. Jede einzelne ergibt einen Unterraum, und der \textbf{Durchschnitt von Untervektorr\"aumen ist wieder ein Unterraum}.

Aber wir k\"onnen es auch direkt beweisen:

\textbf{Beweis: $W$ ist ein Unterraum.}

\textbf{(U1)} $0(1) = 0$ und $0(2) = 0$ \checkmark.

\textbf{(U2)} Seien $p, q \in W$. Dann:
\[
(p+q)(1) = p(1) + q(1) = 0 + 0 = 0, \quad (p+q)(2) = p(2) + q(2) = 0 + 0 = 0 \quad \checkmark
\]

\textbf{(U3)} Sei $p \in W$, $\lambda \in \mathbb{R}$.
\[
(\lambda p)(1) = \lambda \cdot 0 = 0, \quad (\lambda p)(2) = \lambda \cdot 0 = 0 \quad \checkmark
\]

$\Rightarrow$ \textbf{$W$ ist ein Unterraum von $\mathbb{R}[x]$.} \hfill $\blacksquare$

%-----------------------------------------------------------------------------
\subsection*{(c) \quad $W = \{p(x) \in \mathbb{R}[x] : p(1) \cdot p(2) = 0\}$}

\textbf{Was bedeutet die Bedingung?}\\
Das \textbf{Produkt} von $p(1)$ und $p(2)$ muss $0$ sein. Das hei\ss t: $p(1) = 0$ \textbf{ODER} $p(2) = 0$ (oder beides). Mindestens einer der zwei Werte muss $0$ sein.

\textbf{Achtung:} Das ist ein gro\ss er Unterschied zu (b)! In (b) mussten BEIDE $0$ sein (UND). Hier muss nur EINER $0$ sein (ODER). ODER-Bedingungen sind fast nie Untervektorr\"aume!

\begin{merkbox}
Produkt-Bedingung ($p(1) \cdot p(2) = 0$) ist \textbf{nichtlinear} $\Rightarrow$ Addition geht wahrscheinlich schief!
\end{merkbox}

\textbf{$W$ ist KEIN Unterraum.}

\textbf{Gegenbeispiel (U2 verletzt):}
\[
p(x) = x - 1, \quad q(x) = x - 2
\]

$p \in W$? $\quad p(1) \cdot p(2) = 0 \cdot 1 = 0$ \checkmark{} \quad (weil $p(1) = 0$)

$q \in W$? $\quad q(1) \cdot q(2) = (-1) \cdot 0 = 0$ \checkmark{} \quad (weil $q(2) = 0$)

Jetzt:
\[
(p + q)(x) = (x - 1) + (x - 2) = 2x - 3
\]

\[
(p+q)(1) = 2 \cdot 1 - 3 = -1, \quad (p+q)(2) = 2 \cdot 2 - 3 = 1
\]

\[
(p+q)(1) \cdot (p+q)(2) = (-1) \cdot 1 = -1 \neq 0 \quad \Lightning
\]

\textbf{Warum passiert das?} $p$ hat Nullstelle bei $1$, $q$ hat Nullstelle bei $2$. Aber die Summe $p + q$ hat Nullstelle bei $x = 3/2$ -- weder bei $1$ noch bei $2$!

$\Rightarrow$ \textbf{$W$ ist kein Unterraum.} \hfill $\blacksquare$

%-----------------------------------------------------------------------------
\subsection*{(d) \quad $W = \{p(x) \in \mathbb{R}[x] : p(1) \leq p(2)\}$}

\textbf{Was bedeutet die Bedingung?}\\
Der Wert des Polynoms an der Stelle $1$ muss kleiner oder gleich dem Wert an der Stelle $2$ sein.

\begin{merkbox}
Ungleichung $\Rightarrow$ Skalarmultiplikation mit negativer Zahl dreht die Richtung um!
\end{merkbox}

\textbf{$W$ ist KEIN Unterraum.}

\textbf{Gegenbeispiel (U3 verletzt):} Nimm $p(x) = x$ und $\lambda = -1$.

$p \in W$? $\quad p(1) = 1 \leq 2 = p(2)$ \checkmark.

$(-1) \cdot p(x) = -x$. Dann: $(-p)(1) = -1$ und $(-p)(2) = -2$.

$-1 \leq -2$? \textbf{NEIN!} ($-1 > -2$) \quad \Lightning

\textbf{Warum?} Multiplizieren mit $-1$ dreht die Ungleichung um: aus $1 \leq 2$ wird $-1 \geq -2$.

$\Rightarrow$ \textbf{$W$ ist kein Unterraum.} \hfill $\blacksquare$

%-----------------------------------------------------------------------------
\subsection*{(e) \quad $W = \{p(x) \in \mathbb{R}[x] : p(1) = p(2) + p(3)\}$}

\textbf{Was bedeutet die Bedingung?}\\
Der Wert bei $x = 1$ muss gleich der Summe der Werte bei $x = 2$ und $x = 3$ sein. Umgeschrieben: $p(1) - p(2) - p(3) = 0$. Das ist eine \textbf{lineare Bedingung} an $p$ (weil Einsetzen linear ist), und rechts steht $0$.

\textbf{Beweis: $W$ ist ein Unterraum.}

\textbf{(U1)} Nullpolynom: $0 = 0 + 0$ \checkmark.

\textbf{(U2)} Seien $p, q \in W$, also $p(1) = p(2) + p(3)$ und $q(1) = q(2) + q(3)$.
\begin{align*}
(p+q)(1) &= p(1) + q(1) = \bigl(p(2) + p(3)\bigr) + \bigl(q(2) + q(3)\bigr) \\
&= \bigl(p(2) + q(2)\bigr) + \bigl(p(3) + q(3)\bigr) = (p+q)(2) + (p+q)(3) \quad \checkmark
\end{align*}

\textbf{(U3)} Sei $p \in W$, $\lambda \in \mathbb{R}$.
\[
(\lambda p)(1) = \lambda \cdot p(1) = \lambda \bigl(p(2) + p(3)\bigr) = \lambda p(2) + \lambda p(3) = (\lambda p)(2) + (\lambda p)(3) \quad \checkmark
\]

$\Rightarrow$ \textbf{$W$ ist ein Unterraum von $\mathbb{R}[x]$.} \hfill $\blacksquare$

%-----------------------------------------------------------------------------
\subsection*{(f) \quad $W = \{p(x) \in \mathbb{R}[x] : p'(1) = p(0)\}$}

\textbf{Was bedeutet die Bedingung?}\\
$p'(x)$ ist die \textbf{Ableitung} von $p(x)$. Die Bedingung sagt: Die Ableitung ausgewertet bei $x = 1$ muss gleich dem Wert des Polynoms bei $x = 0$ sein.

Beispiel: $p(x) = x^2$. Dann $p'(x) = 2x$, also $p'(1) = 2$ und $p(0) = 0$. Ist $2 = 0$? Nein $\Rightarrow$ $p \notin W$.

Anderes Beispiel: $p(x) = x$. Dann $p'(x) = 1$, also $p'(1) = 1$ und $p(0) = 0$. Ist $1 = 0$? Nein $\Rightarrow$ $p \notin W$.

Beispiel das klappt: $p(x) = x^2 - x$. Dann $p'(x) = 2x - 1$, also $p'(1) = 1$ und $p(0) = 0$. Hmm, $1 \neq 0$. Probiere: $p(x) = 0$: $p'(1) = 0 = p(0)$ \checkmark.

\begin{merkbox}
$p'(1) = p(0)$ ist \"aquivalent zu $p'(1) - p(0) = 0$. Sowohl Ableiten als auch Einsetzen sind \textbf{lineare Operationen}. Also ist die Bedingung linear und homogen $\Rightarrow$ Unterraum!
\end{merkbox}

\textbf{Beweis: $W$ ist ein Unterraum.}

\textbf{(U1)} Nullpolynom: $0' = 0$, also $0'(1) = 0 = 0(0)$ \checkmark.

\textbf{(U2)} Seien $p, q \in W$, also $p'(1) = p(0)$ und $q'(1) = q(0)$.

Die Ableitung ist linear: $(p+q)' = p' + q'$. Also:
\[
(p+q)'(1) = p'(1) + q'(1) = p(0) + q(0) = (p+q)(0) \quad \checkmark
\]

\textbf{(U3)} Sei $p \in W$, $\lambda \in \mathbb{R}$. Es gilt $(\lambda p)' = \lambda p'$, also:
\[
(\lambda p)'(1) = \lambda \cdot p'(1) = \lambda \cdot p(0) = (\lambda p)(0) \quad \checkmark
\]

$\Rightarrow$ \textbf{$W$ ist ein Unterraum von $\mathbb{R}[x]$.} \hfill $\blacksquare$

\newpage

%=============================================================================
%  TEIL 5: AUFGABE 13 (Funktionen-Unterräume)
%=============================================================================
\part*{Teil 5: Aufgabe 13 -- Untervektorr\"aume in $\mathrm{FUNC}(\mathbb{R})$}

Hier ist $V = \mathrm{FUNC}(\mathbb{R})$, der Vektorraum \textbf{aller Funktionen} $f \colon \mathbb{R} \to \mathbb{R}$. Die ``Vektoren'' sind Funktionen! Der Nullvektor ist die \textbf{Nullfunktion} $f(x) = 0$ f\"ur alle $x$.

Addition: $(f + g)(x) = f(x) + g(x)$ (punktweise addieren).\\
Skalarmultiplikation: $(\lambda f)(x) = \lambda \cdot f(x)$.

%-----------------------------------------------------------------------------
\subsection*{(a) \quad $W = \{f \in V : f(-x) = -f(x) \text{ f\"ur alle } x \in \mathbb{R}\}$ (ungerade Funktionen)}

\textbf{Was bedeutet die Bedingung?}\\
Eine Funktion hei\ss t \textbf{ungerade}, wenn sie ``punktsymmetrisch zum Ursprung'' ist. Wenn man $x$ durch $-x$ ersetzt, dreht sich das Vorzeichen des Funktionswerts um.

Beispiele:
\begin{itemize}
    \item $f(x) = x$ ist ungerade: $f(-x) = -x = -f(x)$ \checkmark
    \item $f(x) = x^3$ ist ungerade: $(-x)^3 = -x^3$ \checkmark
    \item $f(x) = \sin(x)$ ist ungerade: $\sin(-x) = -\sin(x)$ \checkmark
    \item $f(x) = x^2$ ist NICHT ungerade: $(-x)^2 = x^2 \neq -x^2$
\end{itemize}

\begin{merkbox}
$f(-x) = -f(x)$ ist eine \textbf{lineare Bedingung} an $f$ (auf beiden Seiten steht was Lineares in $f$). Also: Unterraum!
\end{merkbox}

\textbf{Beweis: $W$ ist ein Unterraum.}

\textbf{(U1)} Nullfunktion: $0(-x) = 0 = -0(-x)$? Nein, genauer: $f(x) = 0$ f\"ur alle $x$. Dann $f(-x) = 0 = -0 = -f(x)$ \checkmark.

\textbf{(U2)} Seien $f, g \in W$, also $f(-x) = -f(x)$ und $g(-x) = -g(x)$ f\"ur alle $x$.
\[
(f+g)(-x) = f(-x) + g(-x) = -f(x) + (-g(x)) = -(f(x) + g(x)) = -(f+g)(x) \quad \checkmark
\]

\textbf{(U3)} Sei $f \in W$, $\lambda \in \mathbb{R}$.
\[
(\lambda f)(-x) = \lambda \cdot f(-x) = \lambda \cdot (-f(x)) = -(\lambda f)(x) \quad \checkmark
\]

$\Rightarrow$ \textbf{$W$ ist ein Unterraum von $\mathrm{FUNC}(\mathbb{R})$.} \hfill $\blacksquare$

%-----------------------------------------------------------------------------
\subsection*{(b) \quad $W = \{f \in V : f(x) \neq 0 \text{ f\"ur alle } x \neq 0\}$}

\textbf{Was bedeutet die Bedingung?}\\
Die Funktion darf \textbf{nur} an der Stelle $x = 0$ den Wert $0$ haben. An allen anderen Stellen muss der Funktionswert $\neq 0$ sein.

\textbf{$W$ ist KEIN Unterraum.}

\textbf{Gegenbeispiel 1 (U1 verletzt):} Die Nullfunktion $f(x) = 0$ hat $f(1) = 0$, obwohl $1 \neq 0$. Also erf\"ullt die Nullfunktion die Bedingung \textbf{nicht}: $f \notin W$.

Schon damit ist klar: $W$ ist kein Unterraum.

\textbf{Gegenbeispiel 2 (U2 verletzt, zur Illustration):} Selbst wenn wir (U1) ignorieren:
\[
f(x) = x \in W \quad \text{(denn $f(x) = x \neq 0$ f\"ur $x \neq 0$ \checkmark)}
\]
\[
g(x) = -x \in W \quad \text{(denn $g(x) = -x \neq 0$ f\"ur $x \neq 0$ \checkmark)}
\]
\[
(f + g)(x) = x + (-x) = 0 \quad \text{f\"ur ALLE } x
\]

Also $(f+g)(1) = 0$, obwohl $1 \neq 0$. Die Summe ist die Nullfunktion, die die Bedingung nicht erf\"ullt! \Lightning

$\Rightarrow$ \textbf{$W$ ist kein Unterraum.} \hfill $\blacksquare$

%-----------------------------------------------------------------------------
\subsection*{(c) \quad $W = \{f \in V : f(x) = f(x+1) \text{ f\"ur alle } x \in \mathbb{R}\}$}

\textbf{Was bedeutet die Bedingung?}\\
Die Funktion ist \textbf{periodisch mit Periode 1}. Das hei\ss t: Wenn man $x$ um $1$ nach rechts verschiebt, \"andert sich der Funktionswert nicht. Der Graph wiederholt sich alle $1$ Einheiten.

Beispiele:
\begin{itemize}
    \item $f(x) = \sin(2\pi x)$ ist periodisch mit Periode 1 \checkmark
    \item Jede konstante Funktion $f(x) = c$ \checkmark{} (denn $c = c$)
    \item $f(x) = x$ ist NICHT periodisch ($x \neq x + 1$)
\end{itemize}

\begin{merkbox}
$f(x) = f(x+1)$ ist eine lineare Bedingung an $f$ $\Rightarrow$ Unterraum!
\end{merkbox}

\textbf{Beweis: $W$ ist ein Unterraum.}

\textbf{(U1)} Nullfunktion: $0 = 0$ f\"ur alle $x$ \checkmark{} (die Nullfunktion ist trivialerweise periodisch).

\textbf{(U2)} Seien $f, g \in W$, also $f(x) = f(x+1)$ und $g(x) = g(x+1)$ f\"ur alle $x$.
\[
(f+g)(x) = f(x) + g(x) = f(x+1) + g(x+1) = (f+g)(x+1) \quad \checkmark
\]

\textbf{(U3)} Sei $f \in W$, $\lambda \in \mathbb{R}$.
\[
(\lambda f)(x) = \lambda \cdot f(x) = \lambda \cdot f(x+1) = (\lambda f)(x+1) \quad \checkmark
\]

$\Rightarrow$ \textbf{$W$ ist ein Unterraum von $\mathrm{FUNC}(\mathbb{R})$.} \hfill $\blacksquare$

%-----------------------------------------------------------------------------
\subsection*{(d) \quad $W = \{f \in V : f(x) + 1 = f(x+1) \text{ f\"ur alle } x \in \mathbb{R}\}$}

\textbf{Was bedeutet die Bedingung?}\\
Wenn man $x$ um $1$ erh\"oht, erh\"oht sich der Funktionswert um genau $1$. Beispiel: $f(x) = x$ erf\"ullt das: $x + 1 = (x + 1)$.

\textbf{Aber Achtung:} Umgeschrieben: $f(x+1) - f(x) = 1$. Die rechte Seite ist $1 \neq 0$! Das ist eine \textbf{inhomogene} Bedingung.

\textbf{$W$ ist KEIN Unterraum.}

\textbf{Gegenbeispiel (U1 verletzt):} Teste die Nullfunktion $f(x) = 0$:
\[
f(x) + 1 = 0 + 1 = 1, \quad f(x + 1) = 0
\]

$1 = 0$? \textbf{NEIN!} \quad \Lightning

Die Nullfunktion erf\"ullt die Bedingung nicht, weil $0 + 1 = 1 \neq 0$.

$\Rightarrow$ \textbf{$W$ ist kein Unterraum.} \hfill $\blacksquare$

\newpage

%=============================================================================
%  ZUSAMMENFASSUNG
%=============================================================================
\part*{Zusammenfassung}

\begin{merkbox}[title={\"Ubersicht aller Ergebnisse}]
\textbf{Aufgabe 10 ($\mathbb{R}^n$):}\\
(a) JA \quad (b) NEIN \quad (c) NEIN \quad (d) NEIN \quad (e) NEIN \quad (f) NEIN \quad (g) JA \quad (h) NEIN \quad (i) NEIN

\medskip
\textbf{Aufgabe 11 ($\mathbb{R}^{n \times n}$):}\\
(a) JA \quad (b) NEIN \quad (c) JA \quad (d) NEIN \quad (e) JA \quad (f) JA

\medskip
\textbf{Aufgabe 12 ($\mathbb{R}[x]$):}\\
(a) JA \quad (b) JA \quad (c) NEIN \quad (d) NEIN \quad (e) JA \quad (f) JA

\medskip
\textbf{Aufgabe 13 ($\mathrm{FUNC}(\mathbb{R})$):}\\
(a) JA \quad (b) NEIN \quad (c) JA \quad (d) NEIN
\end{merkbox}

\begin{merkbox}[title={Die goldenen Regeln f\"ur Unterraum-Aufgaben}]
\textbf{Sofort JA (Unterraum), wenn:}
\begin{itemize}
    \item Homogene lineare Gleichungen (nur Variablen-Terme, rechts steht $0$)
    \item Kern einer linearen Abbildung ($Av = 0$, $\mathrm{Spur}(A) = 0$, $p(1) = 0$, etc.)
    \item Lineare Eigenschaft von Funktionen ($f(-x) = -f(x)$, $f(x) = f(x+1)$)
\end{itemize}

\medskip
\textbf{Sofort NEIN (kein Unterraum), wenn:}
\begin{itemize}
    \item Inhomogene Bedingung (Konstante $\neq 0$ auf einer Seite): Nullvektor meistens nicht drin
    \item Ungleichung ($\leq$, $<$): Skalarmultiplikation mit $-1$ dreht Richtung um
    \item Nichtlineare Bedingung (Produkte, Quadrate, Betr\"age): Addition zerst\"ort die Eigenschaft
    \item $\neq$-Bedingung: Addition kann die ``verbotene'' Gleichheit erzeugen
    \item Verschiedene Typen ($\mathbb{R}^2 \subseteq \mathbb{R}^3$): Keine echte Teilmenge
\end{itemize}
\end{merkbox}

\end{document}
